\section{Conclusion}

In this work, we introduced \mind, the first dataset for developing and evaluating generalist agents for the web. 
We also proposed \model, an agent that leverages the power of (large) language models for effectively tackling this task. 
Our work opens up a wide range of promising future directions, including integrating multi-modal information, reinforcement learning with feedback from real websites, and specialized LMs for web understanding and action taking.
We hope that \mind\ will serve as a valuable platform for the research community to advance towards generalist agents for the web.

\section*{Acknowledgements}
The authors would thank colleagues from the OSU NLP group for constructive feedback and all contributors from the Amazon Mechanical Turk platform who participated in our study and assisted in data collection. 
This research was sponsored in part by NSF OAC 2112606, NSF CAREER \#1942980, ARL W911NF2220144 and Ohio Supercomputer Center~\cite{OhioSupercomputerCenter1987}. The views and conclusions contained herein are those of the authors and should not be interpreted as representing the official policies,
either expressed or implied, of the U.S. government. The U.S. Government is authorized to reproduce and distribute reprints for Government purposes notwithstanding any copyright notice herein.